%==========================================================================
\section{Introduction}
\label{sec:intro}
%==========================================================================

Quasi-periodic oscillations (QPOs) are narrow peaks in the X-ray power spectra of accreting compact objects\footnote{A precise definition, together with the timing and spectral properties of these features, can be found in \citet{ingram}.}.
They were first found in neutron-star and black-hole (BH) X-ray binaries (XRBs), where they are common, and have since been reported in a small number of active galactic nuclei (AGNs).

AGNs are usually regarded as the supermassive counterparts of BH-XRBs: in both classes a BH accretes through a disc, and the relevant physics is expected to be the same.
In particular, the dependence of the Kerr metric on the BH mass can be absorbed by choosing gravitational units ($R_g \doteq GM/c^2$, $t_g \doteq R_g/c$), so that the description of motion becomes independent of $M$.

However, the hidden mass scaling is important for observations:
a vast group of QPO models binds the measured frequency to combinations of Kerr epyciclic frequencies, whose explicit mass dependece is $M^{-1}$.
They therefore predict that a feature seen at frequency $\nu_\textup{XRB}$ around a BH of $10\,M_\odot$ appears around an AGN of mass $M$ at $\nu_0=\nu_\textup{XRB}\,(10\,M_\odot/M)$, i.e. in the $10^{-7}$-$10^{-3}$~Hz range for $M\sim10^{5}$-$10^{8}\,M_\odot$.

The same situation holds for other characteristic timescales of the flow, since they are also obtained from Kerr frequencies and thus also scale as $M$ once the disc parameters (aspect ratio, viscosity, surface-density profile) are fixed.
What does not scale is the comparison with timescales external to the flow: the duration of an X-ray observation, the time over which misaligned gas is supplied to the inner disc, and the lifetime of the active phase.

%--------------------------------------------------------------------------
\subsection{Observations}
Table~\ref{tab:catalog} collects the X-ray QPOs reported in AGNs so far, with the redshift of the host and the published BH mass estimates.
The first robust detection, and still the best studied, is the $\sim1$~h QPO of RE~J1034+396 \citep{Gierlinski2008, Alston2014}.
Similar signals at $\nu_\textup{obs}\sim10^{-4}$~Hz were later reported in MS~2254.9-3712 \citep{Alston2015}, Mrk~766 \citep{Zhang2017}, MCG-06-30-15 \citep{Gupta2018}, 1H~0707-495 \citep{Pan2016,Zhang2018a}, ESO~113-G010 \citep{Zhang2020}, NGC~1365 \citep{Yan2024}, NGC~5506 \citep{Zhang2025}, Mrk~142 \citep{Zhong2024} and NGC~4151 \citep{Yan2025}.
Longer periods, from hours to weeks, were found in XMMU~J134736.6+173403 \citep[J1347 hereafter;][]{Carpano2018}, 2XMM~J123103.2+110648 \citep[J1231;][]{Lin2013}, NGC~4945 \citep{Smith2020} and 2MASX~J12571076+2724177 \citep[J1257;][]{j1257}.
At the other end, the changing-look AGN 1ES~1927+654 shows a millihertz QPO whose frequency drifted from 0.9 to 2.4~mHz within two years and then stabilised \citep{Masterson2025,Masterson2026}.
QPOs have also been claimed in tidal disruption events and in extreme-UV light curves of AGNs; we do not consider them here and refer to the compilation of \citet{Smith2021}.

Two features of this sample are relevant for what follows.
First, the reported significances are often computed only over the portion of the light curve where the QPO is present, and systematic re-analyses of archival data confirm only the QPO of RE~J1034+396 \citep{GonzalezMartin2012,Zhang2023}.
A re-analysis based on Gaussian processes also lowers the significance of the QPOs of NGC~1365 and NGC~4945 to about $1.7\sigma$ and $2.5\sigma$, respectively \citep{Gurpide2025}.
Second, all detections come from long archival observations taken for other purposes, which favours bright, highly variable narrow-line Seyfert~1 galaxies \citep{Smith2021}.
The sample is therefore small, heterogeneous and affected by selection, and the mass estimates of individual sources often disagree by an order of magnitude or more (Table~\ref{tab:catalog}).

\begin{table*}
\caption{X-ray QPOs reported in AGNs.}
\label{tab:catalog}
\centering
\begin{tabular}{lcccll}
\hline\hline
Source & $z$ & $\nu_\textup{obs}$ [Hz] & $M$ [$10^6\,M_\odot$] & QPO ref. & Mass ref. \\
\hline
RE J1034+396           & 0.0424  & $2.7\times10^{-4}$                    & 0.44--55   & 1, 2   & 3 \\
MS 2254.9-3712         & 0.039   & $1.5\times10^{-4}$                    & 4--10      & 4      & 5, 6 \\
Mrk 766                & 0.0129  & $1.55\times10^{-4}$                   & 1.3--7.4   & 7      & 8, 9, 10 \\
MCG-06-30-15           & 0.00775 & $2.7\times10^{-4}$                    & 1.2--58    & 11     & 12, 13 \\
ESO 113-G010           & 0.0257  & $1.24\times10^{-4}$, $6.8\times10^{-5}$ & 4.1--10  & 14     & 15 \\
J1347  & 0.045   & $1.17\times10^{-5}$, $3.9\times10^{-6}$ & 6.2--28  & 16     & 16 \\
J1231 & 0.119   & $7.3\times10^{-5}$                    & 0.03--0.1  & 17     & 18, 19 \\
NGC 4945               & 0.0019  & $2.74\times10^{-7}$                   & 0.7--7.4  & 20     & 21, 22, 23 \\
J1257                  & 0.0207  & $3.3\times10^{-5}$                    & 1.0--5.9  & 24     & 25 \\
1ES 1927+654           & 0.0194  & $0.9$--$2.4\times10^{-3}$             & 0.65--5.1 & 26, 27 & 28 \\
1H 0707-495            & 0.0411  & $2.63\times10^{-4}$, $1.2\times10^{-4}$ & 2--86    & 29, 30 & 31, 32, 33 \\
NGC 1365               & 0.0055  & $2.19\times10^{-4}$                   & 2--63      & 34     & 35, 36, 37 \\
NGC 5506               & 0.0062  & $6.2\times10^{-5}$                    & 4--100     & 38     & 39 \\
Mrk 142                & 0.045   & $7.0\times10^{-5}$                    & 0.6--3.9   & 40     & 41, 42 \\
NGC 4151               & 0.0033  & $5.9\times10^{-4}$                    & 2.5--51    & 43     & 44, 45, 46 \\
\hline
\end{tabular}
\tablefoot{
$\nu_\textup{obs}$ is the observed frequency; two values denote two QPOs in the same source, a range a drifting QPO.
The mass range covers all published independent estimates, including discordant ones; estimates derived from the QPO itself are excluded.
References:
(1)~\citet{Gierlinski2008}; (2)~\citet{Alston2014}; (3)~\citet{Czerny2016};
(4)~\citet{Alston2015}; (5)~\citet{Grupe2010}; (6)~\citet{Shields2003};
(7)~\citet{Zhang2017}; (8)~\citet{BentzKatz2015}; (9)~\citet{Bentz2009}; (10)~\citet{Giacche2014};
(11)~\citet{Gupta2018}; (12)~\citet{Bentz2016}; (13)~\citet{Das2025};
(14)~\citet{Zhang2020}; (15)~\citet{Cackett2013};
(16)~\citet{Carpano2018};
(17)~\citet{Lin2013}; (18)~\citet{Ho2012}; (19)~\citet{Cao2025};
(20)~\citet{Smith2020}; (21)~\citet{Greenhill1997}; (22)~\citet{Braatz2003}; (23)~\citet{Gaspar2022};
(24)~\citet{j1257}; (25)~\citet{massa};
(26)~\citet{Masterson2025}; (27)~\citet{Masterson2026}; (28)~\citet{Li2022};
(29)~\citet{Pan2016}; (30)~\citet{Zhang2018a}; (31)~\citet{ZhouWang2005}; (32)~\citet{DoneJin2016}; (33)~\citet{Titarchuk2025};
(34)~\citet{Yan2024}; (35)~\citet{Fazeli2019}; (36)~\citet{Combes2019}; (37)~\citet{Risaliti2009};
(38)~\citet{Zhang2025}; (39)~\citet{Gofford2015};
(40)~\citet{Zhong2024}; (41)~\citet{Du2014}; (42)~\citet{Li2018};
(43)~\citet{Yan2025}; (44)~\citet{Bentz2006}; (45)~\citet{Onken2014}; (46)~\citet{Roberts2021}.
}
\end{table*}

%--------------------------------------------------------------------------
\subsection{The model}
Many models have been proposed for QPOs in XRBs.
We adopt the one most widely used for type-C QPOs, the Lense-Thirring precession of a geometrically thick inner flow misaligned with the BH spin \citep{PIF}.
We remain agnostic on whether it is the correct description of AGN QPOs: we use it because its predictions follow from a few transparent assumptions and because it contains explicit conditions for the precession to be observable.
According to \citet{Motta2017}, rigid precession at the observed frequency requires that bending waves cross the flow within one period and that the flow does not align with the BH equator within one period.

%--------------------------------------------------------------------------
\subsection{The question}
Within this model the scaling with mass is exact: once the disc parameters are fixed, the observed frequency, the wave-crossing time and the alignment time depend on $M$ and $\nu_0$ only through the product $P=M\nu_0$ (Sect.~\ref{sec2:scaling}).
A configuration that produces a QPO around a $10\,M_\odot$ BH has therefore an exact counterpart around any AGN with the same $P$, and the conditions for the QPO to exist are satisfied, or violated, in the same way.
If QPOs are common in XRBs, why are so few observed in AGNs?

We consider three classes of explanations, which are not mutually exclusive.
\begin{enumerate}
    \item \textit{Restricted parameter space.}
    For given disc parameters, only part of the $(M,\nu_0)$ plane admits a valid precessing configuration, i.e. one where bending waves cross the flow and the flow does not align within one precession period.
    If the allowed region excludes the masses and frequencies where most AGNs lie, QPOs are rare intrinsically.
    \item \textit{Alignment.}
    The flow aligns with the BH spin on the timescale $t_\textup{align}$.
    Even where both conditions above hold, $t_\textup{align}$ must also be compared with the time between successive episodes that bring misaligned gas to the inner disc and with the duration of the active phase.
    In XRBs the companion continuously supplies gas \citep{Maccarone2002, Martin2008}, whereas in AGNs a misaligned inner flow may form only occasionally, for instance after a tidal disruption event or a change in the fuelling; the fraction of time during which a QPO is visible would then be of order $t_\textup{align}$ times the rate of such episodes.
    The QPO of 1ES~1927+654, which appeared after a major change of the accretion flow \citep{Masterson2025}, would fit this picture.
    \item \textit{Detectability.}
    At fixed $P$ the period grows linearly with $M$, so for massive AGNs a QPO may exist but last too few cycles within an observation to be detected \citep{VaughanUttley2005}.
\end{enumerate}

\subsection{This work}
We first use the model to constrain the BH spin and the extent of the precessing flow for each source of Table~\ref{tab:catalog}.
We then map the valid region of the $(M,\nu_0)$ plane, identify which condition fails where no configuration is valid, and test how the results depend on the poorly known parameters of the flow.
Finally, we compare the allowed region with the observational window and estimate the alignment timescales, in order to assess the three explanations above.

The paper is organised as follows.
Section~\ref{sec:model} describes the model and its exact scalings.
Section~\ref{sec:results} presents the constraints on the individual sources, while Section~\ref{sec:general} generalises the analysis to the complete mass-frequency space.
Section~\ref{sec:discussion} finally discusses the three explanations and summarises our conclusions.
Appendix~\ref{app:choices} discusses modelisation choices compared to previous works, and Appendix~\ref{app:params} tests the soundness of our results against variations of the disc's parameters.